\subsection{Logarithmic divisors with a small-curvature residual}
\label{subsec:small-curvature-intrinsic-traces}

An ambient residual can be invisible to every multiplier ideal even
when its restriction to $Y$ is identically $-\infty$. The restricted
ideals then recover the imposed logarithmic divisor directly.

\begin{proposition}\label{prop:small-curvature-intrinsic-traces}\index{Trace!small-curvature residual}\index{Multiplier ideal!normal crossing divisor}
Let $q\in\QPSH(X)$ satisfy \eqref{eq:small-curvature-type}, and let
$\Delta$ and $\sigma_\Delta$ be as in
\cref{thm:small-curvature-strict-type}. For every $t>0$,
\[
 \mathcal I\bigl(t(\sigma_\Delta+q)\bigr)
                  =\mathcal I(t\sigma_\Delta).
\]
If $\Delta$ has simple normal crossings and $Y$ is contained in none
of its positive-coefficient components, then for every divisor $E$
on every smooth connected compact modification $\rho:Z\to Y$,
\[
 \tau_{Y,\sigma_\Delta+q}(E)
     =\sum_ja_j\operatorname{ord}_E\rho^*(D_j|_Y).
\]
The formula imposes no condition on the naive restriction $q|_Y$.
\end{proposition}
\begin{proof}
The first equality follows from
\cref{lem:small-curvature-zero-lelong,lem:zero-lelong-ideal-neutrality}.
In normal crossing coordinates, the weight for $t\sigma_\Delta$
is a smooth positive factor times $\prod_j|z_j|^{-2ta_j}$.
One-variable polar integration gives the exact formula
\[
 \mathcal I(t\sigma_\Delta)
       =\mathcal O_X\left(-\sum_j\lfloor ta_j\rfloor D_j\right).
\]
Indeed, integrability of a germ of transverse order $k$ requires
$k+1>ta_j$, so $k\geq\lfloor ta_j\rfloor$; equality at the preceding
order is divergent when $ta_j$ is an integer. Necessity follows on an
open set of the other variables where the leading coefficient is nonzero.
After division by these powers, the remaining product weight is
integrable, proving sufficiency.

Since $Y$ is contained in no $D_j$, each restricted defining equation
is a nonzero holomorphic germ on $Y$. Its pullback to $Z$ is nonzero
and has a finite nonnegative order $N_{j,E}$ along $E$.
Restriction of the displayed principal ideal and then pullback therefore
give the exact order $\sum_j\lfloor ta_j\rfloor N_{j,E}$.
Dividing by $t$ proves the trace formula, including along all real weights.
Moreover $\nu(\sigma_\Delta+q,Y)=0$: outside the finitely many divisors
both summands have zero point Lelong number. Thus these orders are the
intrinsic data of \cref{thm:intrinsic-trace-profile}. When $\Delta=0$,
the ideal is $\mathcal O_X$ and every trace is zero.
\end{proof}
